\section*{第 \ref{ch:g12:radioactive-decay} 章——放射性衰变}

\begin{solution}{exo:g12:radioactive-decay:1}
$T_{1/2} = \ln 2/\lambda = 0.693/\num{1.0e-6} = \qty{6.9e5}{s}
\approx 8.0$ 天——碘-131 的“时钟”。
\end{solution}

\begin{solution}{exo:g12:radioactive-decay:2}
$T_{1/2} = 3.8 \times \num{86400} = \qty{3.28e5}{s}$；
$\lambda = 0.693/\num{3.28e5} = \qty{2.1e-6}{s^{-1}}$；
$\mathcal A = \lambda N = \num{2.1e-6} \times \num{1.0e10} \approx
\qty{2.1e4}{Bq}$。
\end{solution}

\begin{solution}{exo:g12:radioactive-decay:3}
$\dd N/\dd t = N_0(-\lambda)e^{-\lambda t} = -\lambda N$ ---
方程成立；$N(0) = N_0 e^0 = N_0$；
$N(T_{1/2})/N_0 = e^{-\lambda \ln 2/\lambda} = e^{-\ln 2} =
\tfrac12$。
\end{solution}

\begin{solution}{exo:g12:radioactive-decay:4}
$90$ 年 $= 3\,T_{1/2}$：$\mathcal A = \num{4.0e5}/2^3 =
\qty{5.0e4}{Bq}$。$300$ 年 $= 10\,T_{1/2}$：
$\mathcal A = \num{4.0e5}/1024 \approx \qty{3.9e2}{Bq}$。
\end{solution}

\begin{solution}{exo:g12:radioactive-decay:5}
斜率 $= (18.85 - 20.00)/10 = -0.115$，因此
$\lambda = \qty{0.115}{h^{-1}} = 0.115/3600 = \qty{3.2e-5}{s^{-1}}$；
$T_{1/2} = 0.693/0.115 = 6.0$ 小时——锝-99m。
\end{solution}

\begin{solution}{exo:g12:radioactive-decay:6}
$N = (1.0/226) \times \num{6.02e23} = \num{2.7e21}$；
$T_{1/2} = 1600 \times \num{3.156e7} = \qty{5.05e10}{s}$，因此
$\lambda = \qty{1.37e-11}{s^{-1}}$；
$\mathcal A = \lambda N \approx \qty{3.7e10}{Bq}$——一居里，即
玛丽·居里所用一克镭的放射性活度。
\end{solution}

\begin{solution}{exo:g12:radioactive-decay:7}
两者都预期剩下一半：$10$，以及 \num{1.0e20}。计数的涨落
$\sim 1/\sqrt{n}$：对 $10$ 次衰变，约为 $30\%$——
$7$ 或 $13$ 并不出人意料；对 \num{1.0e20}，约为
$10^{-10}$。只有大样本才能在
可测量的精度上服从该规律。
\end{solution}

\begin{solution}{exo:g12:radioactive-decay:8}
$m = 140 \times \num{1.2e-4} = \qty{1.7e-2}{g}$；
$N = (0.0168/40) \times \num{6.02e23} = \num{2.5e20}$。
$\lambda = 0.693/(\num{1.25e9} \times \num{3.156e7}) =
\qty{1.76e-17}{s^{-1}}$；$\mathcal A = \lambda N \approx
\qty{4.4e3}{Bq}$——约为人体总放射性活度的一半（其余大部分
来自碳-14）。
\end{solution}

\begin{solution}{exo:g12:radioactive-decay:9}
$\lambda = 0.693/6.0 = \qty{0.116}{h^{-1}}$。在 $15$ 小时时：
$\mathcal A = 500\,e^{-0.116 \times 15} = 500 \times 0.177 \approx
\qty{88}{MBq}$。低于 \qty{1}{MBq} 时：
$t = \ln(500)/0.116 \approx \qty{54}{hours}$——约为 $9$
个半衰期。
\end{solution}

\begin{solution}{exo:g12:radioactive-decay:10}
木炭：$3.4/13.6 = \tfrac14$，恰为 $2$ 个半衰期：
$t = \num{1.15e4}$ 年。骨骼：
$t = \num{5730} \times \ln(1/0.30)/\ln 2 = \num{5730} \times 1.74
\approx \num{1.0e4}$ 年。
\end{solution}

\begin{solution}{exo:g12:radioactive-decay:11}
$\lambda = 0.693/3.8 = \qty{0.182}{d^{-1}}$；
$t = \ln(100)/\lambda = 4.61/0.182 \approx 25$ 天（约 $6.6$
个半衰期）。但周围花岗岩和土壤中的铀衰变链会
持续产生氡：一次性通风会在
数周内被抵消，因此通风必须长期进行。
\end{solution}

\begin{solution}{exo:g12:radioactive-decay:12}
由 $N/N_0 = e^{-\lambda t}$ 得 $\lambda t = \ln(N_0/N)$，而
由 $\lambda = \ln 2/T_{1/2}$ 得 $t = T_{1/2}\ln(N_0/N)/\ln 2$。此处
$t = \num{5730} \times \ln(20)/\ln 2 = \num{5730} \times 4.32
\approx \num{2.5e4}$ 年。
\end{solution}

\begin{solution}{exo:g12:radioactive-decay:13}
$N_{\mathrm K} = N_0 e^{-\lambda t}$ 且
$N_{\mathrm{Ar}} = N_0 - N_{\mathrm K}$，因此
$N_{\mathrm{Ar}}/N_{\mathrm K} = e^{\lambda t} - 1$。比值为 $3.0$：
$e^{\lambda t} = 4$，故 $\lambda t = 2\ln 2$，
$t = 2\,T_{1/2} = \num{2.5e9}$ 年。
\end{solution}

\begin{solution}{exo:g12:radioactive-decay:14}
$\ln \mathcal A = 5.99$、$5.42$、$4.84$、$4.26$、$3.69$：每隔
\qty{60}{s} 下降 $0.57$--$0.58$——近似一致，呈直线。斜率：
$\lambda = 2.30/240 = \qty{9.6e-3}{s^{-1}}$；
$T_{1/2} = 0.693/\num{9.6e-3} \approx \qty{72}{s}$。
\end{solution}

\begin{solution}{exo:g12:radioactive-decay:15}
(a) $2^{-10} \approx \num{1e-3}$。(b) $(1/2)^8 = 1/256 \approx
0.4\%$。(c) 预期为 $1000/8 = 125$——但该规律对于
哪些原子会衰变并不作答：它精确预言群体的行为，却完全无法预言个体。
\end{solution}

\begin{solution}{pb:g12:radioactive-decay:1}
\textbf{1.} 与大气的交换（进食、呼吸）使
该比例始终维持在大气中的数值；死亡后摄入停止，
衰变便起主导作用。

\textbf{2.} $T_{1/2} = \num{5730} \times \num{3.156e7} =
\qty{1.81e11}{s}$，故 $\lambda = 0.693/\num{1.81e11} =
\qty{3.8e-12}{s^{-1}}$.

\textbf{3.} $\dd N/\dd t = -\lambda N_0 e^{-\lambda t} = -\lambda N$,
且 $N(0) = N_0$。

\textbf{4.} $\mathcal A_0 = 13.6/60 = \qty{0.227}{Bq}$ 每克；
$N = \mathcal A_0/\lambda = 0.227/\num{3.8e-12} \approx \num{5.9e10}$
个原子每克。

\textbf{5.} $(1/12) \times \num{6.02e23} = \num{5.0e22}$ 个碳
原子；比例 $\num{5.9e10}/\num{5.0e22} \approx \num{1.2e-12}$
——即每 $10^{12}$ 个中有一个，与所述相符。

\textbf{6.} 由 $\mathcal A = \mathcal A_0 e^{-\lambda t}$ 得
$\lambda t = \ln(\mathcal A_0/\mathcal A)$，再由 $\lambda =
\ln 2/T_{1/2}$ 代入，得
$t = T_{1/2}\ln(\mathcal A_0/\mathcal A)/\ln 2$.

\textbf{7.} $t = \num{5730} \times \ln(13.6/9.1)/\ln 2 =
\num{5730} \times 0.580 \approx \num{3.3e3}$ 年。

\textbf{8.} 木材的“死亡”——即树木被砍伐之时，而非
雕刻之时；由古老木材雕成的雕像（或以古木雕成的赝品）
其年代早于或晚于其原材料的年代。

\textbf{9.} $t = \num{5730} \times \ln(13.6/10.4)/\ln 2 =
\num{5730} \times 0.387 \approx \num{2.2e3}$ 年——完全
与二十二个世纪前的书吏相符。

\textbf{10.} $\Delta t \approx \Delta\mathcal A/(\lambda\mathcal A)$
其中 $\lambda = \ln 2/5730 = \num{1.21e-4}$ 每年：
$\Delta t = (0.2/10.4)/\num{1.21e-4} \approx \num{1.6e2}$ 年——
即年代精确到一两个世纪之内。

\textbf{11.} $t = \num{5730} \times \ln(13.6/13.4)/\ln 2 \approx
120$ 年——在 $\pm 160$ 年
的不确定度之内，与现代物品无法区分：羊皮纸是近代的，所谓“古物”乃是赝品。

\textbf{12.} $13.6/2^n < 0.2$ 要求 $2^n > 68$：$n = 7$ 个半衰期
（因 $2^7 = 128$），约 $\num{4.0e4}$ 年。

\textbf{13.} 十个半衰期：约 \num{5.7e4} 年——
即实际可测的极限，约六万年。

\textbf{14.} $\num{6.6e7}/\num{5730} \approx \num{11500}$ 个半衰期。
一克中的 \num{5.9e10} 个原子在
$2^n > \num{5.9e10}$ 之后耗尽，即 $n \approx 36$ 个半衰期
（$\approx \num{2e5}$ 年）：远在恐龙时代之前，已无一个
碳-14 原子留存——无从计数。

\textbf{15.} 一种慢得多的“时钟”：钾-40 或铀-238，其
半衰期长达数十亿年。

\textbf{16.} 熔融的熔岩使氩气逸出，因此刚凝固的
晶体含有钾而氩为零：该比值
在凝固时刻即为 $0$——时钟被归零。

\textbf{17.} 如练习所示：$N_{\mathrm K} = N_0 e^{-\lambda t}$，
$N_{\mathrm{Ar}} = N_0 - N_{\mathrm K}$，故
$N_{\mathrm{Ar}}/N_{\mathrm K} = e^{\lambda t} - 1$.

\textbf{18.} $e^{\lambda t} = 1.60$:
$t = \num{1.25e9} \times \ln(1.60)/\ln 2 = \num{1.25e9} \times 0.678
\approx \num{8.5e8}$ 年。

\textbf{19.} 不行，有两个原因：在 \num{8.5e8} 年时，碳-14 已历经
$10^5$ 个半衰期——早已沉寂；在 \num{3.3e3} 年时，钾-40 仅历经
\num{2.6e-6} 个半衰期——氩比值
$\approx \lambda t \sim \num{2e-6}$ 小到无法测量。每只时钟
仅能读取约十分之一个半衰期至约十个半衰期的范围。

\textbf{20.} 雕像 $\approx \num{3.3e3}$ 年，莎草纸
$\approx \num{2.2e3}$ 年，岩石 $\approx \num{8.5e8}$ 年——
三只时钟，一条规律：$N = N_0\,e^{-\lambda t}$，逆向读取。
\end{solution}
